\documentclass[11pt,aspectratio=43]{beamer}

% A simple academic Beamer layout with a fresh teal-and-gold palette.
\usepackage[T1]{fontenc}
\usepackage{lmodern}
\usepackage{amsmath}
\usepackage{xcolor}

\definecolor{Teal}{RGB}{0,112,105}
\definecolor{DeepTeal}{RGB}{8,67,68}
\definecolor{Gold}{RGB}{214,157,56}
\definecolor{SoftTeal}{RGB}{231,244,242}

\setbeamertemplate{navigation symbols}{}
\setbeamertemplate{items}[circle]
\setbeamercolor{normal text}{fg=black,bg=white}
\setbeamercolor{structure}{fg=Teal}
\setbeamercolor{frametitle}{fg=white,bg=Teal}
\setbeamercolor{title}{fg=white,bg=Teal}
\setbeamercolor{item}{fg=Gold}
\setbeamercolor{subitem}{fg=Teal}
\setbeamercolor{footline}{fg=white,bg=DeepTeal}
\setbeamerfont{frametitle}{size=\large,series=\bfseries}
\setbeamerfont{title}{size=\Large,series=\bfseries}

% Wide colored title band, matching the reference slide layout.
\setbeamertemplate{frametitle}{%
  \nointerlineskip%
  \begin{beamercolorbox}[wd=\paperwidth,ht=2.7ex,dp=1.1ex,leftskip=0.35cm,rightskip=0.25cm]{frametitle}%
    \insertframetitle\strut%
  \end{beamercolorbox}%
}

% Compact three-part footer with author/ID, deck title, and slide count.
\setbeamertemplate{footline}{%
  \leavevmode%
  \hbox{%
    \begin{beamercolorbox}[wd=.38\paperwidth,ht=2.6ex,dp=1.1ex,leftskip=.25cm]{footline}%
      \scriptsize Sagar Jaiswal  |  2025UCP1944
    \end{beamercolorbox}%
    \begin{beamercolorbox}[wd=.40\paperwidth,ht=2.6ex,dp=1.1ex,center]{footline}%
      \scriptsize Introduction to Beamer
    \end{beamercolorbox}%
    \begin{beamercolorbox}[wd=.22\paperwidth,ht=2.6ex,dp=1.1ex,rightskip=.25cm plus1fil]{footline}%
      \scriptsize \insertframenumber/\inserttotalframenumber
    \end{beamercolorbox}%
  }%
}

\title{Introduction to Beamer}
\author{Sagar Jaiswal}
\institute{ID: 2025UCP1944}
\date{\today}

\begin{document}

% Slide 1: Title slide - centered title panel and presenter details.
\begin{frame}
  \vspace{1.35cm}
  \begin{center}
  \begin{beamercolorbox}[wd=.92\textwidth,rounded=true,shadow=true,sep=1.0em,center]{title}
    {\usebeamerfont{title}\inserttitle}
  \end{beamercolorbox}
  \end{center}
  \vspace{0.65cm}
  \begin{center}
    {\Large \insertauthor}\par
    \vspace{0.28cm}
    {\normalsize \insertinstitute}\par
    \vspace{0.45cm}
    {\normalsize \insertdate}
  \end{center}
\end{frame}

% Slide 2: Introduction
\begin{frame}{Introduction}
  \vspace{0.55cm}
  \textbf{What is Beamer?}
  \begin{itemize}
    \item Beamer is a LaTeX document class used to create presentations.
    \item It allows us to create structured slides.
    \item It supports equations, tables, figures and more.
  \end{itemize}
\end{frame}

% Slide 3: Document and page styles
\begin{frame}{Document and Page Styles}
  \vspace{0.22cm}
  \textbf{Document Class}
  \begin{itemize}
    \item Defines the overall structure of a document.
    \item Examples:
    \begin{itemize}
      \item article
      \item report
      \item book
    \end{itemize}
  \end{itemize}

  \vspace{0.18cm}
  \textbf{Page Style}
  \begin{itemize}
    \item Controls the appearance of headers and footers.
    \item Examples: plain, empty, headings and fancy.
  \end{itemize}
\end{frame}

% Slide 4: Mathematical equations
\begin{frame}{Mathematical Equations}
  \vspace{0.35cm}
  Beamer also supports mathematical expressions.

  \begin{center}
    \vspace{0.15cm}
    \[
      E = mc^2
    \]
    \vspace{0.08cm}
    \[
      \sum_{i=1}^{n} i = \frac{n(n+1)}{2}
    \]
  \end{center}

  \vspace{0.05cm}
  \textbf{Advantages}
  \begin{itemize}
    \item Mathematical expressions are easy to format.
    \item Equations can be displayed clearly on slides.
  \end{itemize}
\end{frame}


% Additional slide: Blocks and emphasis
\begin{frame}{Using Blocks to Highlight Information}
  Beamer provides blocks that separate key information from the rest of the slide.
  \vspace{0.15cm}
  \begin{block}{Regular Block}
    Use this style for definitions, explanations, or important notes.
  \end{block}
  \begin{alertblock}{Important Notice}
    Use an alert block sparingly when a point needs special attention.
  \end{alertblock}
  \begin{exampleblock}{Example}
    A short example helps connect a concept with a practical use.
  \end{exampleblock}
\end{frame}

% Additional slide: Columns
\begin{frame}{Arranging Content in Columns}
  Columns place related material side by side instead of in one long list.
  \vspace{0.18cm}
  \begin{columns}[T,onlytextwidth]
    \column{0.48\textwidth}
      \begin{block}{Text and Lists}
        \begin{itemize}
          \item Explain an idea.
          \item Give supporting points.
          \item Keep each point concise.
        \end{itemize}
      \end{block}
    \column{0.48\textwidth}
      \begin{block}{Examples and Visuals}
        \begin{itemize}
          \item Show a small table.
          \item Add a diagram.
          \item Include an equation.
        \end{itemize}
      \end{block}
  \end{columns}
  \vspace{0.10cm}
  \small Use columns when two pieces of information are easier to understand together.
\end{frame}

% Additional slide: Overlays
\begin{frame}{Revealing Content Step by Step}
  Overlays let a presenter reveal selected items gradually instead of showing everything at once.
  \vspace{0.12cm}
  \begin{block}{A list with step-by-step appearance}
    \begin{itemize}
      \item<1-> First, introduce the main idea.
      \item<2-> Next, explain a supporting detail.
      \item<3-> Finally, show the conclusion.
    \end{itemize}
  \end{block}
  \textbf{Useful overlay commands}
  \begin{itemize}
    \item \texttt{\textbackslash pause} pauses before showing the next content.
    \item \texttt{\textbackslash only<2->\{...\}} displays material from overlay 2 onward.
    \item \texttt{\textbackslash visible<2->\{...\}} reserves space before material appears.
  \end{itemize}
\end{frame}

% Additional slide: Themes
\begin{frame}{Themes and Visual Customization}
  Beamer themes control the overall appearance of a presentation, including colors, navigation elements, and slide headings.
  \vspace{0.12cm}
  \begin{block}{Selecting a theme}
    \small
    \texttt{\textbackslash usetheme\{Madrid\}}\par
    \texttt{\textbackslash usecolortheme\{seahorse\}}
  \end{block}
  \textbf{Design suggestions}
  \begin{itemize}
    \item Choose a theme that suits the purpose of the presentation.
    \item Use a small set of matching colors.
    \item Keep headings and body text visually distinct.
    \item Avoid putting too much information on one slide.
  \end{itemize}
\end{frame}

% Additional slide: Presentation tips
\begin{frame}{Tips for an Effective Presentation}
  A good slide deck supports the speaker rather than replacing the speaker.
  \vspace{0.12cm}
  \begin{itemize}
    \item \textbf{Keep slides focused:} present one main idea on each slide.
    \item \textbf{Use readable text:} choose font sizes visible from a distance.
    \item \textbf{Prefer concise points:} explain details verbally instead of filling slides with paragraphs.
    \item \textbf{Use visuals carefully:} add diagrams, tables, or images when they clarify the topic.
    \item \textbf{Check consistency:} keep colors, headings, and spacing uniform.
    \item \textbf{Rehearse:} practise moving through the slides before presenting.
  \end{itemize}
\end{frame}

% Slide 5: Conclusion
\begin{frame}{Conclusion}
  \begin{center}
    \vspace{0.2cm}
    {\LARGE\bfseries Thank You!}\par
    \vspace{0.35cm}
    {\large Questions?}
  \end{center}

  \vspace{0.45cm}
  \textbf{Key Takeaways}
  \begin{itemize}
    \item Beamer is used to create presentations using LaTeX.
    \item Each slide is created using the frame environment.
    \item We can include text, equations, tables and figures.
  \end{itemize}
\end{frame}

\end{document}
